\documentclass[11pt]{article}
\usepackage[margin=1in]{geometry}
\usepackage{amsmath,amssymb,amsthm}
\usepackage[hidelinks]{hyperref}
\newtheorem{theorem}{Theorem}
\newtheorem{lemma}{Lemma}
\begin{document}
\begin{center}
{\Large The coefficient at the penultimate Markov-sail vertex}\\[0.5em]
A proof of MathDB 366808\\
September 12, 2026
\end{center}

\paragraph{Task 1.}
Prove the Location of 4 conjecture of Evans, Veselov, and Winn
\cite[Conjecture 8.4]{EVW}. We use their numerator convention
\[
 M_{a/b}(x,y,z)=
 \frac{P_{a/b}(x^2,y^2,z^2)}{x^{a-1}y^{b-1}z^{a+b-1}},
 \qquad
 P_{a/b}(u,v,w)=\sum_{i,j} A_{ij}u^iv^jw^{a+b-1-i-j}.
\]
For coprime $2\leq a<b$, the conjecture specifies a coefficient in the
interior of the critical triangle
$i<a$, $j<b$, $bi+aj>ab$.
If $b/a=[h_1,\ldots,h_n]$, with convergents $p_k/q_k$, its sail
coordinates are
\[
 A_m^{\mathrm{sail}}=(q_{2m-1},b-p_{2m-1}),\qquad
 B_m^{\mathrm{sail}}=(a-q_{2m},p_{2m}).
\]
The claimed coefficient is $4$ at $A_m^{\mathrm{sail}}$ when $n=2m$,
and at $B_m^{\mathrm{sail}}$ when $n=2m+1$.
The source defines sail coefficients only in that interior; for
$a=1$ the critical triangle is empty, as it explicitly notes in
Section 9.

\begin{theorem}
For every coprime pair $2\leq a<b$, there is a unique pair of integers
$0<i<a$, $0<j<b$ satisfying
\[
 bi+aj=ab+1.
\]
The coefficient $A_{ij}$ of $P_{a/b}$ at this pair is $4$.
Consequently Conjecture 8.4 holds in both parity cases.
\end{theorem}

\begin{proof}
We first recall two established facts from
\cite[Theorems 3.2 and 4.2]{EVW}. For coprime positive $r,s$, the
Newton polygon of $P_{r/s}(u,v,1)$ is
\[
 \left\{(i,j):i,j\geq0,\quad \frac{i}{r}+\frac{j}{s}\geq1,
                    \quad i+j\leq r+s-1\right\}.
\]
Its coefficients at $(r,0)$ and $(0,s)$ are both $1$.
No assertion of saturation of the polygon is needed.

Work in formal Laurent series in $t$, and set
\[
 \kappa(t)=\frac{1+t^{2a}+t^{2b}}{t^{a+b}},\qquad
 F_{r/s}(t)=\kappa(t)M_{r/s}(t^b,t^a,1).
\]
The numerator convention gives the useful identity
\[
 F_{r/s}(t)
 =(1+t^{2a}+t^{2b})\,t^{-(br+as)}
                     P_{r/s}(t^{2b},t^{2a},1).
 \tag{1}
\]
If $|br-as|=1$, the linear weight $bi+aj$ has a unique minimum
$\min(br,as)$ on the Newton polygon, at one of its two lower
endpoints. The minimum coefficient is $1$. All weight differences
are integers, so (1) implies
\[
 F_{r/s}(t)=t^{-1}+O(t).
 \tag{2}
\]
Here and below the $O$ notation denotes divisibility in the ring
of formal Laurent series.

Let $r/s$ and $u/v$ be the Farey parents of $a/b$, so that
$r+u=a$, $s+v=b$, and $rv-su=\pm1$.
Both parents have positive numerator and denominator since $a\geq2$.
Then $|br-as|=|bu-av|=1$.
The three associated Markov polynomials form a triple satisfying
the generalized Markov equation. After multiplication by $\kappa$,
the series
\[
 U=F_{r/s},\qquad V=F_{u/v},\qquad W=F_{a/b}
\]
therefore satisfy
\[
 U^2+V^2+W^2=UVW,\qquad
 UV(W-2)=(U-V)^2+W^2.
 \tag{3}
\]

For the target polygon, the minimum weight is $ab$ along the segment
from $(a,0)$ to $(0,b)$. Coprimality implies that its only lattice
points are its endpoints, each with coefficient $1$. Thus
$W=2+O(t^2)$. By (2), $U-V=O(t)$ and
$UV=t^{-2}(1+O(t^2))$. Substituting these facts into (3) yields
\[
 W-2=\frac{(U-V)^2+W^2}{UV}
     =4t^2+O(t^4).
 \tag{4}
\]
Because $a,b\geq2$, the factor $1+t^{2a}+t^{2b}$ in (1)
does not change the coefficient of $t^2$ in $W$.
That coefficient is therefore
\[
 \sum_{bi+aj=ab+1} A_{ij}=4.
 \tag{5}
\]

There is exactly one nonnegative integer solution to
$bi+aj=ab+1$. Indeed choose the unique $i\in\{1,\ldots,a-1\}$
with $bi\equiv1\pmod a$ and put $j=(ab+1-bi)/a$.
Then $0<j<b$, and any other integer solution differs by
$(ka,-kb)$. Neither $k>0$ nor $k<0$ preserves nonnegativity.
This point also satisfies $i+j\leq a+b-2$, so it lies in the
critical triangle and in the stated Newton polygon.
Equation (5) proves $A_{ij}=4$.

Finally put $(p,q)=(p_{n-1},q_{n-1})$ for the penultimate convergent.
With the source's indexing $p_0=1,q_0=0$, the continued-fraction
determinant identity gives
\[
 bq-ap=(-1)^n.
\]
If $n$ is even, the point $(q,b-p)=A_m^{\mathrm{sail}}$ has
weight $bq+a(b-p)=ab+1$.
If $n$ is odd, the point $(a-q,p)=B_m^{\mathrm{sail}}$ has
weight $b(a-q)+ap=ab+1$.
Both are interior points, and uniqueness identifies each with the
point in (5). This proves both assertions of the conjecture.
\end{proof}

\paragraph{Task 2.}
Determine explicit formulas for the coefficients at lattice heights
$bi+aj-ab=2,3,\ldots$ by expanding the same formal Markov identity
to higher orders.

\paragraph{Task 3.}
Prove the Factor 4 conjecture for every coefficient in the critical
triangle, extending the height-one result beyond the penultimate
vertex.

\paragraph{Task 4.}
Use valuations adapted to other sail vertices to prove the
Markov Sail Duality identities for arithmetic progressions of
coefficients along sail segments.

\paragraph{Task 5.}
Give a direct perfect-matching interpretation of the four objects
counted by the coefficient at the penultimate sail vertex.

\begin{thebibliography}{9}
\raggedright
\bibitem{EVW}
S. J. Evans, A. P. Veselov, and B. Winn,
\emph{Arithmetic and geometry of Markov polynomials},
arXiv:2501.14882v3 (2025); current 28-page journal-hosted manuscript,
Theorems 3.2 and 4.2, and Conjecture 8.4.
\url{https://armj.math.stonybrook.edu/Articles/250708-Winn/article.pdf}.
\end{thebibliography}
\end{document}
